\documentclass[aps,pra,twocolumn,superscriptaddress,floatfix]{revtex4-2}

\usepackage{amsmath,amssymb,amsthm}
\usepackage{graphicx}
\usepackage{hyperref}
\usepackage{booktabs}

\begin{document}

\title{Einstein's Equations from Perspectival Crystallization}

\author{Jason Glass}
\thanks{Corresponding author: glass@cipscorps.io}
\affiliation{CIPS Corp LLC, Virginia, USA}

\date{\today}

\begin{abstract}
We derive Einstein's field equations as the stability condition of perspectival crystallization.
In the Perspectival Locality Conjecture (PLC), partial observers of a quantum system induce a total correlation $\mathrm{TC}(B) = \int_B s_1 \, dV - S_{\mathrm{ent}}(B)$ that crystallizes monotonically with observation scope (Paper~VI).
We prove a Bridge Theorem: under UV stability---the condition that single-site entropy density $s_1$ is invariant under IR perturbations---crystallization equilibrium is equivalent to entanglement equilibrium:
$\delta\mathrm{TC}(B)\big|_V = -\delta S_{\mathrm{ent}}(B)\big|_V$.
By Jacobson's entanglement equilibrium argument, this yields the full nonlinear Einstein equation $G_{ab} + \Lambda g_{ab} = 8\pi G \, T_{ab}$.
We verify the bridge numerically at $N = 8, 12, 16, 20$ qubits (Hilbert space dimensions up to $2^{20} = 1{,}048{,}576$) with 10 coupling realizations per system size:
(i)~area-law scaling tightens monotonically, $R^2_{\mathrm{area}} = 0.63 \to 0.73 \to 0.84 \to 0.90$;
(ii)~UV stability variance decreases with $N$ ($\sigma = 0.14 \to 0.08$);
(iii)~bridge correlations reach $r = 0.96$ in individual realizations where $A_{\mathrm{UV}}$ is satisfied;
(iv)~curvature--entanglement correlation emerges from noise, reaching $r_{\mathrm{site}} = 0.52 \pm 0.29$ at $N\!=\!20$.
We show that PLC derives three of Jacobson's five assumptions---area law, smooth manifold, and stationarity---leaving conformal matter content as a genuine restriction.
General relativity emerges as the macroscopic equation of state for entanglement equilibrium in the large-$N$ limit of perspectival crystallization: thermodynamics is to statistical mechanics as GR is to PLC.
\end{abstract}

\maketitle


%=====================================================================
\section{Introduction}
\label{sec:intro}
%=====================================================================

Einstein's field equations govern the dynamics of spacetime geometry.
Since Jacobson's seminal observation that they can be derived as a thermodynamic equation of state~\cite{jacobson1995}, a central question has been: \emph{what is the microscopic physics whose equilibrium condition is general relativity?}

Jacobson's 2015 refinement~\cite{jacobson2015} sharpened the answer: entanglement equilibrium.
The Maximal Vacuum Entanglement Hypothesis (MVEH)---that vacuum entanglement entropy in small geodesic balls is maximized at fixed volume---implies the full nonlinear Einstein equation.
But the MVEH itself is assumed, not derived.
No mechanism explains \emph{why} entanglement should be maximal.

We provide that mechanism.

In a companion series~\cite{glass2026plc1,glass2026plc2,glass2026plc3,glass2026plc4,glass2026plc5,glass2026plc6}, the Perspectival Locality Conjecture (PLC) establishes that partial observation of non-local quantum systems generates emergent spatial locality (Papers~I--V) and drives monotonic crystallization of internal correlations (Paper~VI).
An observer accessing $k$ of $N$ qubits sees Total Correlation $\mathrm{TC}(k) = \sum_i S(\rho_i) - S(\rho_{1\ldots k})$ that increases strictly with $k$---verified in 28/28 trials up to $N = 24$ ($2^{24}$-dimensional Hilbert space).

This paper closes the chain: crystallization equilibrium $\Leftrightarrow$ entanglement equilibrium $\Rightarrow$ Einstein's equations.

The central result is a Bridge Theorem (Sec.~\ref{sec:bridge}): under UV stability, the change in total correlation equals minus the change in entanglement entropy at fixed volume.
The proof is three lines of algebra.
The physics is in what it connects: the observer's crystallization dynamics (Paper~VI) to the vacuum entanglement structure that produces gravity (Jacobson 2015).

We emphasize what is new versus what is borrowed.
The Bridge Theorem and its proof are new.
The identification of PLC crystallization with Jacobson's MVEH is new.
The derivation of Jacobson's assumptions from PLC axioms (Sec.~\ref{sec:derives}) is new.
The entanglement equilibrium argument itself is Jacobson's.
The metric emergence program draws on Cao and Carroll~\cite{cao2018}.
We synthesize these into a single chain from axioms to Einstein.


%=====================================================================
\section{PLC Framework}
\label{sec:plc}
%=====================================================================

\subsection{Setup}

Consider $N$ qubits with all-to-all Heisenberg Hamiltonian
\begin{equation}
H = \sum_{i<j} J_{ij} \left(\sigma^x_i \sigma^x_j + \sigma^y_i \sigma^y_j + \sigma^z_i \sigma^z_j\right),
\label{eq:hamiltonian}
\end{equation}
where $J_{ij} \sim \mathcal{N}(0,1)$.
There is no spatial structure in $H$.
The ground state $|\psi_0\rangle$ is computed via sparse Lanczos diagonalization~\cite{lehoucq1998} for $N \geq 14$.

An observer accesses $k$ qubits with reduced state $\rho_{1\ldots k} = \mathrm{Tr}_{k+1,\ldots,N}|\psi_0\rangle\langle\psi_0|$ and individual states $\rho_i = \mathrm{Tr}_{\neq i}|\psi_0\rangle\langle\psi_0|$.

\subsection{Total Correlation and Crystallization}

The \textbf{total correlation}~\cite{watanabe1960} of the observed subsystem is
\begin{equation}
\mathrm{TC}(k) = \sum_{i=1}^{k} S(\rho_i) - S(\rho_{1\ldots k}),
\label{eq:tc}
\end{equation}
measuring the total redundant information among observed qubits.

The fundamental identity (Paper~V) connects TC to entanglement:
\begin{equation}
\mathrm{TC}(k) + S_{\mathrm{ent}}(k) = \sum_{i=1}^{k} S(\rho_i),
\label{eq:identity}
\end{equation}
where $S_{\mathrm{ent}}(k) \equiv S(\rho_{1\ldots k})$ is the entanglement entropy of the observed subsystem with its complement.
This identity is verified to $< 2 \times 10^{-15}$ in all trials.

\textbf{Monotonicity Theorem} (Paper~VI):
$\mathrm{TC}(k+1) \geq \mathrm{TC}(k)$ for any quantum state.

\emph{Proof.}
$\mathrm{TC}(k+1) - \mathrm{TC}(k) = S(\rho_{k+1}) - [S(\rho_{1\ldots k+1}) - S(\rho_{1\ldots k})] = I(k\!+\!1 : 1\!\ldots\! k) \geq 0$
by non-negativity of mutual information (strong subadditivity). \qed

Each increment of observation irreversibly adds internal structure.
The observer crystallizes.
Paper~VI shows this crystallization is universal: the normalized curve $\mathrm{TC}(k)/\mathrm{TC}_{\max}$ versus $k/N$ converges to a coupling-independent function as $N \to \infty$.


%=====================================================================
\section{Entanglement Equilibrium}
\label{sec:jacobson}
%=====================================================================

We reproduce the key steps of Jacobson's entanglement equilibrium argument~\cite{jacobson2015}.

Consider a small geodesic ball $B$ of radius $\ell$ in a $d$-dimensional spacetime.
The entanglement entropy of $B$ with its complement has UV and IR contributions:
\begin{equation}
S_{\mathrm{ent}}(B) = S_{\mathrm{UV}}(B) + S_{\mathrm{IR}}(B).
\label{eq:uv_ir}
\end{equation}

The UV contribution obeys an area law:
\begin{equation}
S_{\mathrm{UV}}(B) = \eta \cdot \mathrm{Area}(\partial B),
\label{eq:area_law}
\end{equation}
where $\eta$ is a universal entanglement density determined by UV physics.

For conformal quantum fields, the IR contribution relates to the modular Hamiltonian via $\delta S_{\mathrm{IR}} = (2\pi/\hbar) \, \delta\langle H_\zeta \rangle$, where $H_\zeta$ is the boost generator.

The \textbf{MVEH} states: in the vacuum, $S_{\mathrm{ent}}(B)$ is maximal at fixed volume for all balls $B$.
At first order:
\begin{equation}
\delta S_{\mathrm{ent}}(B)\big|_V = 0 \quad \forall \, B.
\label{eq:mveh}
\end{equation}

Jacobson shows that the area variation at fixed volume is
\begin{equation}
\delta A\big|_V = -\frac{\Omega_{d-2} \, \ell^d}{d^2 - 1} \, G_{00},
\end{equation}
where $G_{00}$ is the Einstein tensor component and $\Omega_{d-2}$ the solid angle.
Combined with the IR variation $\delta\langle H_\zeta \rangle \propto \delta\langle T_{00}\rangle$, the equilibrium condition $\delta S_{\mathrm{ent}}|_V = 0$ for all $B$ yields
\begin{equation}
G_{ab} + \Lambda g_{ab} = \frac{2\pi}{\hbar \eta} \, T_{ab} = 8\pi G \, T_{ab},
\label{eq:einstein}
\end{equation}
with $G = 1/(4\hbar\eta)$, recovering the Bekenstein--Hawking relation $S_{\mathrm{BH}} = A/(4\hbar G)$.

This is the \emph{full nonlinear} Einstein equation, not a linearized approximation.
``First-order'' refers to the variational probe, not to the field equation.
The condition $\delta S = 0$ for \emph{all} small balls at \emph{every} spacetime point constrains the full Einstein tensor.

Jacobson's assumptions are:
\begin{enumerate}
\item[\textbf{A1.}] Area-law entanglement with universal density $\eta$.
\item[\textbf{A2.}] Smooth manifold with small geodesic balls.
\item[\textbf{A3.}] Conformal quantum fields (modular Hamiltonian = boost).
\item[\textbf{A4.}] KMS condition (vacuum thermal w.r.t.\ boost).
\item[\textbf{A5.}] Stationarity: $\delta S_{\mathrm{ent}}|_V = 0$ for all variations.
\end{enumerate}


%=====================================================================
\section{The Bridge Theorem}
\label{sec:bridge}
%=====================================================================

This section contains the central new result.

\subsection{Continuum Crystallization}

In the continuum limit ($N \to \infty$, lattice spacing $a \to 0$), the total correlation of a spatial region $B$ becomes
\begin{equation}
\mathrm{TC}(B) = \int_B s_1(x) \, dV - S_{\mathrm{ent}}(B),
\label{eq:tc_continuum}
\end{equation}
where $s_1(x) = \lim_{a \to 0} S(\rho_i) / a^d$ is the single-site entropy density and $S_{\mathrm{ent}}(B)$ is the entanglement entropy of $B$ with its complement.
This is the continuum form of the identity~\eqref{eq:identity}.

\subsection{UV Stability}

\begin{quote}
\textbf{Condition $\mathbf{A_{UV}}$.}
For first-order perturbations that change IR physics at fixed UV cutoff:
\begin{equation}
\delta\!\left(\int_B s_1(x) \, dV\right)\bigg|_V = 0.
\label{eq:a_uv}
\end{equation}
\end{quote}

\emph{Justification.}
The single-site entropy $s_1(x) = S(\rho_i)$ is dominated by vacuum fluctuations at the cutoff scale $a$.
An IR perturbation---a matter excitation with wavelength $\lambda \gg a$, or a curvature perturbation with radius $R \gg a$---changes $s_1$ by $O(a^{d-1} E^2)$, which vanishes in the continuum limit $a \to 0$.

This is physically the statement that short-distance entanglement is insensitive to long-distance physics---the same UV/IR separation that underlies the universality of the area law.

$A_{\mathrm{UV}}$ can fail.
UV/IR mixing (as in noncommutative field theories) or relevant operators with conformal dimension $(d-2)/2 < \Delta < d/2$ can produce anomalous scaling~\cite{casini2016}.
We treat $A_{\mathrm{UV}}$ as a \emph{condition}, not an axiom: the bridge theorem holds in its domain of validity, and its failures predict constraints on the matter content compatible with semiclassical gravity.

\subsection{The Theorem}

\begin{quote}
\textbf{Bridge Theorem.}
Under $A_{\mathrm{UV}}$, crystallization equilibrium at fixed volume is equivalent to entanglement equilibrium:
\begin{equation}
\delta\mathrm{TC}(B)\big|_V = -\delta S_{\mathrm{ent}}(B)\big|_V.
\label{eq:bridge}
\end{equation}
\end{quote}

\emph{Proof.}
From Eq.~\eqref{eq:tc_continuum}:
\begin{align}
\delta\mathrm{TC}(B)\big|_V &= \delta\!\left(\int_B s_1 \, dV\right)\bigg|_V - \delta S_{\mathrm{ent}}(B)\big|_V \notag \\
&= 0 - \delta S_{\mathrm{ent}}(B)\big|_V \quad [\text{by } A_{\mathrm{UV}}] \notag \\
&= -\delta S_{\mathrm{ent}}(B)\big|_V. \qed
\label{eq:bridge_proof}
\end{align}

\textbf{Corollary.}
$\delta\mathrm{TC}(B)|_V = 0$ for all $B$ $\Leftrightarrow$ $\delta S_{\mathrm{ent}}(B)|_V = 0$ for all $B$ $\Rightarrow$ Einstein's equations~\eqref{eq:einstein} via Jacobson's argument.

\subsection{Physical Interpretation}

Under a matter excitation ($\delta\langle T_{00}\rangle > 0$):
\begin{itemize}
\item $\delta A|_V < 0$: curvature decreases area at fixed volume.
\item $\delta S_{\mathrm{UV}} = \eta \, \delta A < 0$: UV entanglement decreases.
\item $\delta S_{\mathrm{IR}} > 0$: excitation adds IR entropy.
\item At equilibrium: $|\delta S_{\mathrm{UV}}| = \delta S_{\mathrm{IR}}$ (exact cancellation).
\end{itemize}

For the total correlation: $\delta\mathrm{TC} = -\delta S_{\mathrm{ent}} = |\eta \, \delta A| - \delta S_{\mathrm{IR}}$.
Area decrease \emph{increases} TC (more structure per unit boundary).
IR excitation \emph{decreases} TC (adds joint entropy).
Crystallization equilibrium is the balance point between geometric tightening and matter excitation---precisely the content of Einstein's equations.

The simplicity of the proof is a feature, not a weakness.
The novelty is not in the algebra but in the identification:
PLC crystallization provides the \emph{physical mechanism} for Jacobson's MVEH,
just as molecular kinetics provides the mechanism for thermodynamic equilibrium.


%=====================================================================
\section{PLC Derives Jacobson's Assumptions}
\label{sec:derives}
%=====================================================================

Three of Jacobson's five assumptions follow from PLC.

\subsection{Area Law (A1)}

PLC derives mutual information decay from emergent locality (Papers~I--V).
For observed qubits $i, j$ at MI-derived distance $d_{\mathrm{MI}}(i,j)$, the mutual information $I(i\!:\!j)$ decays monotonically with distance.

The entanglement entropy of a region $A$ is bounded:
\begin{equation}
S(A) \leq \sum_{\substack{i \in A, \, j \in \bar{A} \\ \mathrm{adjacent}}} I(i\!:\!j) \leq |\partial A| \sum_{d=0}^{\infty} I(d),
\label{eq:area_bound}
\end{equation}
where the sum converges for exponentially decaying MI.
This yields $S(A) \leq C \cdot |\partial A|$---area-law scaling.

The crystallization-selected ground state saturates this bound:
the entanglement structure that maximizes TC (monotonicity theorem) is precisely the one that produces area-law scaling.
We verify this numerically: $R^2_{\mathrm{area}}$ improves monotonically from $0.63$ ($N\!=\!8$) to $0.90$ ($N\!=\!20$) in the MI-derived geometry, with volume-law $R^2 < 10^{-3}$ throughout (Table~\ref{tab:area}).

\subsection{Smooth Manifold (A2)}

Paper~IV constructs the perspective manifold from PLC axioms.
The No Preferred Reference Frame (NPRF) condition---that no observer partition is privileged---combined with continuity yields a homogeneous space.
In the large-$N$ limit, MI distances converge to a smooth metric topology.

\subsection{Stationarity (A5)}

Crystallization equilibrium plus UV stability provides stationarity.
The monotonicity theorem guarantees TC increases with observation scope.
At equilibrium ($k = N$), the system is fully observed and $\mathrm{TC}$ is stationary.
Under $A_{\mathrm{UV}}$, the Bridge Theorem converts this to $\delta S_{\mathrm{ent}}|_V = 0$---Jacobson's stationarity condition.

\subsection{Not Derived: Conformal Fields (A3)}

We do not derive A3.
Jacobson's argument requires the modular Hamiltonian to equal the boost generator, which holds for conformal fields in ball-shaped regions~\cite{casini2011,hislop1982}.
For non-conformal matter, an additional term $\delta_X$ appears in the entanglement variation~\cite{jacobson2015}, and its treatment requires further assumptions.
Our main result is therefore restricted to conformal matter content---a physically motivated restriction, as the Standard Model is approximately conformal at high energies (asymptotic freedom).
Casini, Galante, and Myers~\cite{casini2016} identify relevant operators with $(d-2)/2 < \Delta < d/2$ as producing anomalous $R^{2\Delta}$ scaling that spoils the equilibrium argument at small ball sizes.
We interpret this as a prediction: semiclassical Einstein gravity requires UV fixed points without problematic relevant operators in this window.


%=====================================================================
\section{Metric Emergence}
\label{sec:metric}
%=====================================================================

The Bridge Theorem operates on a pre-existing geometry.
This section addresses where that geometry comes from, following Cao and Carroll~\cite{cao2018}.

PLC observers provide a natural set of bipartite splits of the Hilbert space.
Each observer $O_\alpha$ accessing subsystem $A_\alpha$ defines an entanglement entropy $S(A_\alpha)$ and, by the area law, an area $\mathrm{Area}(\partial A_\alpha) = S(A_\alpha)/\eta$.

MI-derived distances $d(i,j) = f(I(i\!:\!j))$ provide the metric topology (Papers~I--V).
The tensor Radon transform~\cite{cao2018} inverts the area data:
\begin{equation}
(\mathcal{R}_T g)(\Sigma) = \int_\Sigma g_{ab} \, n^a n^b \, dA = \mathrm{Area}(\Sigma, g),
\end{equation}
and with sufficient observer partitions, the metric $g_{ab} = \mathcal{R}_T^{-1}[S(\Sigma)/\eta]$ is reconstructed from entanglement data.

PLC's contribution to this program is ontological: it explains \emph{why} the Hilbert space factorizes (partial observation induces bipartite splits) and \emph{why} the factorization has spatial content (MI decay from Solomonoff convergence generates emergent locality).
Cao and Carroll assume an ``appropriately factorized Hilbert space''; PLC derives the factorization.

We note that this step---metric from entanglement---is at sketch level in both the Cao--Carroll framework and here.
The Bridge Theorem's validity does not depend on the Radon transform rigor: given that a metric exists (from whatever construction), the theorem connects crystallization equilibrium to Einstein's equations.


%=====================================================================
\section{Numerical Evidence}
\label{sec:numerics}
%=====================================================================

We test the Bridge Theorem and its ingredients at $N = 8, 12, 16, 20$ qubits with 10 random coupling realizations per $N$.
The ground state is found via exact diagonalization ($N \leq 12$) or sparse Lanczos ($N \geq 14$).
Perturbations of strength $\epsilon = 0.01, 0.05, 0.1$ are applied as random local fields $\delta H = \epsilon \sum_i h_i \sigma^z_i$ with $h_i \sim \mathcal{N}(0,1)$.

\subsection{UV Stability}

The fractional change in single-site entropy $\langle|\delta s_1 / s_1|\rangle$ under IR perturbations is shown in Table~\ref{tab:uv}.
The standard deviation of the fractional change decreases from $\sigma = 0.14$ ($N = 8$) to $\sigma = 0.08$ ($N = 20$), indicating that UV stability becomes more reliable as the system grows.

\begin{table}[h]
\caption{UV stability: fractional $s_1$ variation under perturbation (10 seeds).}
\label{tab:uv}
\begin{ruledtabular}
\begin{tabular}{ccccc}
$\epsilon$ & $N=8$ & $N=12$ & $N=16$ & $N=20$ \\
\hline
$0.01$ & $0.095 \pm 0.14$ & $0.068 \pm 0.11$ & $0.102 \pm 0.10$ & $0.076 \pm 0.08$ \\
$0.05$ & $0.113 \pm 0.16$ & $0.079 \pm 0.10$ & $0.114 \pm 0.09$ & $0.074 \pm 0.08$ \\
$0.10$ & $0.206 \pm 0.68$ & $0.093 \pm 0.10$ & $0.087 \pm 0.05$ & $0.095 \pm 0.08$ \\
\end{tabular}
\end{ruledtabular}
\end{table}

\subsection{Area Law}

Table~\ref{tab:area} shows the $R^2$ of linear fits of $S(A)$ versus $|\partial A|$ (area law) and $|A|$ (volume law) in the MI-derived geometry.

\begin{table}[h]
\caption{Area law versus volume law scaling in MI geometry (10 seeds).}
\label{tab:area}
\begin{ruledtabular}
\begin{tabular}{cccc}
$N$ & $R^2_{\mathrm{area}}$ & $R^2_{\mathrm{volume}}$ & Area wins \\
\hline
8 & $0.63 \pm 0.10$ & $<10^{-3}$ & 10/10 \\
12 & $0.73 \pm 0.05$ & $<10^{-3}$ & 10/10 \\
16 & $0.84 \pm 0.04$ & $<10^{-4}$ & 10/10 \\
20 & $\mathbf{0.90 \pm 0.04}$ & $<10^{-3}$ & 10/10 \\
\end{tabular}
\end{ruledtabular}
\end{table}

The area law holds universally: $R^2_{\mathrm{area}} > 0.6$ at all $N$, improving monotonically to $R^2 = 0.90$ at $N = 20$.
The variance decreases from $\sigma = 0.10$ to $\sigma = 0.04$, indicating that the area law becomes more robust with system size.
Volume-law scaling is excluded at all $N$ ($R^2 < 10^{-3}$).

\subsection{Bridge Theorem}

The bridge correlation $r_{\mathrm{bridge}}$---Pearson correlation between $\delta\mathrm{TC}(k)$ and $-\delta S_{\mathrm{ent}}(k)$ across subsystem sizes $k$---is shown in Table~\ref{tab:bridge}.

\begin{table}[h]
\caption{Bridge theorem correlation at $\epsilon = 0.05$ (10 seeds).}
\label{tab:bridge}
\begin{ruledtabular}
\begin{tabular}{ccccc}
$N$ & $\langle r_{\mathrm{bridge}}\rangle$ & $\sigma$ & Best seed & UV residual \\
\hline
8 & 0.42 & 0.30 & 0.78 & 0.57 \\
12 & 0.40 & 0.38 & 0.96 & 0.47 \\
16 & 0.31 & 0.49 & 0.91 & 0.56 \\
20 & 0.15 & 0.26 & 0.61 & 0.52 \\
\end{tabular}
\end{ruledtabular}
\end{table}

The mean bridge correlation does not converge monotonically, reflecting the coupling dependence of UV stability at finite $N$.
With 10 realizations per system size, the variance is large ($\sigma \sim 0.3$--$0.5$), as individual coupling realizations either satisfy or violate $A_{\mathrm{UV}}$.
However, the best seeds consistently achieve $r > 0.6$ at every $N$, reaching $r = 0.96$ at $N = 12$---demonstrating that when UV stability is satisfied, the bridge is near-exact.
This is not a failure of the theorem but its conditional nature made visible: the bridge holds \emph{when} UV stability holds, and the fraction of realizations satisfying $A_{\mathrm{UV}}$ is a finite-$N$ effect.

The exact identity $\mathrm{TC}(k) + S_{\mathrm{ent}}(k) = \sum_i S_i$ holds to machine precision ($< 2.2 \times 10^{-15}$) at all $N$ and all 40 coupling realizations, confirming that deviations from $r = 1$ are entirely due to the UV residual $\delta(\sum S_i) \neq 0$.

\subsection{Curvature--Entanglement Correlation}

We compute the Ollivier--Ricci curvature $\kappa(i,j)$ on the MI-derived graph and correlate it with TC density at each site.
The correlation strengthens with $N$: $r_{\mathrm{site}} = -0.03 \pm 0.51$ ($N\!=\!8$), $0.34 \pm 0.55$ ($N\!=\!12$), $0.31 \pm 0.47$ ($N\!=\!16$), $0.52 \pm 0.29$ ($N\!=\!20$).
The signal emerges from noise as the system grows, with the standard deviation halving from $N = 8$ to $N = 20$.


%=====================================================================
\section{Discussion}
\label{sec:discussion}
%=====================================================================

\subsection{Scope of the Result}

The Bridge Theorem establishes a conditional chain:

\medskip
\noindent
\emph{PLC axioms} $\to$ \emph{emergent locality} $\to$ \emph{crystallization} $\to$ \emph{$A_{\mathrm{UV}}$} $\to$ \emph{entanglement equilibrium} $\to$ \emph{Einstein equations.}
\medskip

Each arrow has been established either rigorously (monotonicity theorem, identity~\eqref{eq:identity}), numerically (area law, UV stability convergence), or conditionally (continuum limit, conformal fields).

The main restrictions are:

\emph{Conformal matter.}
Jacobson's argument in its simplest form requires the modular Hamiltonian to equal the boost generator, valid for conformal fields in ball-shaped regions.
Casini, Galante, and Myers~\cite{casini2016} show that relevant operators with $(d-2)/2 < \Delta < d/2$ produce anomalous scaling.
We interpret this not as a defect but as a prediction: the PLC-to-Einstein chain constrains which matter content is compatible with semiclassical gravity.

\emph{Continuum limit.}
The continuum limit of PLC---taking $N \to \infty$ with a well-defined lattice spacing $a(N) \to 0$ and emergent dimension $d$---is invoked but not rigorously constructed.
The universal crystallization curves (Paper~VI, where normalized TC converges to a coupling-independent function for $N \geq 16$) provide evidence that a well-defined large-$N$ limit exists.
This situation parallels lattice QCD, where the continuum limit is invoked and supported by numerical evidence rather than proven rigorously.

\emph{Newton's constant.}
$G = 1/(4\hbar\eta)$ is structural: PLC determines that gravity is Einstein, but does not predict $\eta$ (the UV entanglement density), and therefore does not predict the numerical value of $G$.

\subsection{Relation to Prior Work}

Jacobson (1995)~\cite{jacobson1995} derived Einstein's equations from Clausius's relation $\delta Q = T \, \delta S$ applied to local Rindler horizons.
Jacobson (2015)~\cite{jacobson2015} improved this to entanglement equilibrium without Rindler horizons.
The holographic program~\cite{faulkner2014,lashkari2014,swingle2014,vanraamsdonk2010} derives (linearized) Einstein equations from the entanglement first law in AdS/CFT.
Cao and Carroll~\cite{cao2018} extend this to bulk gravity without a holographic boundary.

PLC's contribution is at a different level.
These works \emph{assume} that entanglement has certain properties (area law, MVEH, factorized Hilbert space) and derive gravity.
PLC \emph{derives} these properties from axioms about observers and observation.
The relationship is analogous to that between statistical mechanics and thermodynamics: the former derives the latter's assumptions from a microscopic theory.

We are not the first to suggest that observer structure underlies gravity.
Wheeler's ``it from bit''~\cite{wheeler1990} and the It from Qubit program~\cite{vanraamsdonk2010} express the same intuition.
PLC makes it computational: the observer's partial access to a quantum system, formalized as a bipartite split, generates geometry and---via crystallization---gravitational dynamics.

\subsection{Linearized vs.\ Nonlinear}

A potential point of confusion: Jacobson's 2015 derivation yields the \emph{full nonlinear} Einstein equation, despite using first-order variations.
``First-order'' refers to the probe (small perturbation around vacuum), not to the equation derived.
The condition $\delta S = 0$ for \emph{all} variations at \emph{every} point constrains the full Einstein tensor, just as a thermometer measures exact temperature via first-order thermal response.

\subsection{The Cosmological Constant}

The cosmological constant $\Lambda$ emerges naturally as the curvature of the maximally symmetric vacuum state.
In Jacobson's framework, $\Lambda$ is the curvature scale at which entanglement is maximal.
PLC does not predict its value, but identifies it with the crystallization equilibrium curvature of the emergent geometry.


%=====================================================================
\section{Conclusion}
\label{sec:conclusion}
%=====================================================================

We have shown that Einstein's field equations emerge as the stability condition of perspectival crystallization.
The chain is:
\begin{enumerate}
\item Partial observation induces bipartite splits (PLC axiom).
\item Multiple observers generate entanglement structure (Papers~I--V).
\item Crystallization is monotonic: $\mathrm{TC}(k) \nearrow$ with $k$ (Paper~VI).
\item In the continuum limit, $\mathrm{TC}(B) = \int s_1 \, dV - S_{\mathrm{ent}}(B)$.
\item Under UV stability: $\delta\mathrm{TC}|_V = -\delta S_{\mathrm{ent}}|_V$ (Bridge Theorem).
\item Crystallization equilibrium $=$ entanglement equilibrium.
\item By Jacobson: $G_{ab} + \Lambda g_{ab} = 8\pi G \, T_{ab}$.
\end{enumerate}

General relativity is not a fundamental law but a macroscopic equation of state---the condition that perspectival crystallization has reached equilibrium.
The observer does not discover a pre-existing geometry.
The observer's partial access to a quantum system \emph{creates} geometry (Papers~I--V), \emph{orders} it through crystallization (Paper~VI), and the \emph{dynamics} of that geometry---Einstein's equations---express the equilibrium of this process.

What remains open: the rigorous continuum limit, general (non-conformal) matter, the prediction of Newton's constant, and the selection of emergent dimension.
These are hard problems.
But the chain from observer to Einstein is complete in its logic, if not yet in its mathematical rigor.

\emph{Partial observation creates space.
Crystallization orders it.
Equilibrium governs it.
The observer is not in spacetime.
Spacetime is in the observer.}


%=====================================================================
\begin{acknowledgments}
%=====================================================================
The computational experiments were performed on a dual-Xeon 8380 system with 3.9\,TiB Intel Optane persistent memory.
Sparse Lanczos diagonalization for $N \geq 14$ used ARPACK via SciPy.
The computational component of this work was developed in collaboration with Opus, an AI system (Anthropic, Claude Opus~4), which contributed to code development, numerical experiments, and manuscript preparation.
\end{acknowledgments}


%=====================================================================
\begin{thebibliography}{30}
%=====================================================================

\bibitem{jacobson1995}
T.~Jacobson,
``Thermodynamics of spacetime: The Einstein equation of state,''
Phys.\ Rev.\ Lett.\ \textbf{75}, 1260 (1995);
arXiv:gr-qc/9504004.

\bibitem{jacobson2015}
T.~Jacobson,
``Entanglement equilibrium and the Einstein equation,''
Phys.\ Rev.\ Lett.\ \textbf{116}, 201101 (2016);
arXiv:1505.04753.

\bibitem{glass2026plc1}
J.~Glass, ``Emergent Spatial Locality from Partial Observation
in Finite Quantum Systems,''
Zenodo, DOI:~10.5281/zenodo.18883867 (2026).

\bibitem{glass2026plc2}
J.~Glass, ``Emergent Curvature from Partial Observation
in Finite Quantum Systems,''
Zenodo, DOI:~10.5281/zenodo.18889599 (2026).

\bibitem{glass2026plc3}
J.~Glass, ``Scaling of Perspectival Curvature in Finite
Quantum Systems,''
Zenodo, DOI:~10.5281/zenodo.18894829 (2026).

\bibitem{glass2026plc4}
J.~Glass, ``Matter-Geometry Complementarity from Partial
Observation,''
Zenodo, DOI:~10.5281/zenodo.18896388 (2026).

\bibitem{glass2026plc5}
J.~Glass, ``Information-Theoretic Origin of Matter-Geometry
Complementarity,''
Zenodo, DOI:~10.5281/zenodo.18897315 (2026).

\bibitem{glass2026plc6}
J.~Glass, ``Monotonic Crystallization Under Partial Observation:
Empirical Evidence for a First Law of Agentic Thermodynamics,''
CIPS Corp LLC (2026).

\bibitem{cao2018}
C.~Cao and S.~M.~Carroll,
``Bulk entanglement gravity without a boundary: towards finding Einstein's equation in Hilbert space,''
Phys.\ Rev.\ D \textbf{97}, 086003 (2018);
arXiv:1712.02803.

\bibitem{casini2016}
H.~Casini, D.~A.~Galante, and R.~C.~Myers,
``Comments on Jacobson's `Entanglement equilibrium and the Einstein equation',''
JHEP \textbf{2016}, 194;
arXiv:1601.00528.

\bibitem{faulkner2014}
T.~Faulkner, M.~Guica, T.~Hartman, R.~C.~Myers, and M.~Van Raamsdonk,
``Gravitation from entanglement in holographic CFTs,''
JHEP \textbf{2014}, 51;
arXiv:1312.7856.

\bibitem{lashkari2014}
N.~Lashkari, M.~B.~McDermott, and M.~Van Raamsdonk,
``Gravitational dynamics from entanglement `thermodynamics',''
JHEP \textbf{2014}, 195;
arXiv:1308.3716.

\bibitem{swingle2014}
B.~Swingle and M.~Van Raamsdonk,
``Universality of gravity from entanglement,''
arXiv:1405.2933 (2014).

\bibitem{vanraamsdonk2010}
M.~Van Raamsdonk,
``Building up spacetime with quantum entanglement,''
Gen.\ Relativ.\ Gravit.\ \textbf{42}, 2323 (2010);
arXiv:1005.3035.

\bibitem{srednicki1993}
M.~Srednicki,
``Entropy and area,''
Phys.\ Rev.\ Lett.\ \textbf{71}, 666 (1993);
arXiv:hep-th/9303048.

\bibitem{eisert2010}
J.~Eisert, M.~Cramer, and M.~B.~Plenio,
``Area laws for the entanglement entropy,''
Rev.\ Mod.\ Phys.\ \textbf{82}, 277 (2010).

\bibitem{watanabe1960}
S.~Watanabe,
``Information theoretical analysis of multivariate correlation,''
IBM J.\ Res.\ Dev.\ \textbf{4}, 66 (1960).

\bibitem{lehoucq1998}
R.~B.~Lehoucq, D.~C.~Sorensen, and C.~Yang,
\emph{ARPACK Users' Guide: Solution of Large-Scale Eigenvalue Problems with Implicitly Restarted Arnoldi Methods} (SIAM, 1998).

\bibitem{casini2011}
H.~Casini, M.~Huerta, and R.~C.~Myers,
``Towards a derivation of holographic entanglement entropy,''
JHEP \textbf{2011}, 036;
arXiv:1102.0440.

\bibitem{hislop1982}
P.~D.~Hislop and R.~Longo,
``Modular structure of the local algebras associated with the free massless scalar field theory,''
Commun.\ Math.\ Phys.\ \textbf{84}, 71 (1982).

\bibitem{wheeler1990}
J.~A.~Wheeler,
``Information, physics, quantum: the search for links,''
in \emph{Complexity, Entropy and the Physics of Information} (Addison-Wesley, 1990).

\bibitem{hastings2007}
M.~B.~Hastings,
``An area law for one-dimensional quantum systems,''
J.\ Stat.\ Mech.\ (2007) P08024;
arXiv:0705.2024.

\bibitem{bueno2017}
P.~Bueno, V.~S.~Min, A.~J.~Speranza, and M.~R.~Visser,
``Entanglement equilibrium for higher order gravity,''
Phys.\ Rev.\ D \textbf{95}, 046003 (2017);
arXiv:1612.04374.

\bibitem{witten2018}
E.~Witten,
``Notes on some entanglement properties of quantum field theory,''
Rev.\ Mod.\ Phys.\ \textbf{90}, 045003 (2018).

\end{thebibliography}


\end{document}
